% =========================
% 8_appendix.tex --- follows the bibliography.
% =========================
\appendix

\section{Scope of the Closest Systems}
\label{app:entailment}

Table~\ref{tab:coverage} keeps four questions separate: handling acquisition rate, receiving
placement or channel context, accepting run-time activity names, and evaluating multiple target
example counts without parameter updates. A zero-shot result is not marked as a $k$ curve, and a
linear probe or fine-tuning result is not marked as no-update enrollment.

\begin{table*}[t]
  \centering
  \caption{Reported scope of related systems. Entries summarize the cited evaluation rather than
  infer unreported capabilities. ``Task fit'' includes a probe or fine-tuning; ``zero only'' is not
  a target-example curve.}
  \label{tab:coverage}
  \scriptsize
  \setlength{\tabcolsep}{3.5pt}
  \begin{tabularx}{\textwidth}{@{}l
      >{\raggedright\arraybackslash}X
      >{\raggedright\arraybackslash}X
      >{\raggedright\arraybackslash}X
      >{\raggedright\arraybackslash}X@{}}
    \toprule
    System & Rate/input treatment & Configuration context & Label interface & Target adaptation \\
    \midrule
    LiMU-BERT~\cite{limubert} & common 20\,Hz & none & fixed downstream & task fit \\
    CrossHAR~\cite{crosshar} & common-rate inputs & none & fixed downstream & task fit \\
    harnet~\cite{sslwearables} & wrist contract & wrist accelerometer & fixed downstream & task fit \\
    LSM / LSM-2~\cite{lsm1,lsm2} & multimodal, incomplete data & modality availability & fixed tasks & probe/fine-tune \\
    UniMTS~\cite{unimts} & resampled inputs & placement graph & text-conditioned tasks & task fit \\
    NormWear~\cite{normwear} & common-rate input & variable channels & text interface & zero/task evaluation \\
    ZeroHAR~\cite{zerohar} & configured input & categorical context & run-time names & zero only \\
    CHARM / SensorsVoice~\cite{charm,sensorsvoice} & channel-tokenized & channel descriptions & task dependent & task evaluation \\
    MoPFormer~\cite{mopformer} & all data at 100\,Hz & wearable context & fixed tasks & task fit \\
    FlowState~\cite{flowstate} & rate equivariant & sampling rate & forecasting output & not HAR \\
    RAG-HAR~\cite{raghar} & retrieved descriptors & sensor evidence & run-time names & zero only \\
    ZARA~\cite{zara} & multi-sensor evidence & sensor evidence & run-time names & zero only \\
    SensorLM~\cite{sensorlm} & wearable pretraining & sensor/task prompts & task dependent & task adaptation \\
    \midrule
    \textbf{HALO} & stored-grid bands in Hz & per-sensor text + measurements & run-time names & $k=0,1,2,4,8$, no update \\
    \bottomrule
  \end{tabularx}
\end{table*}

\section{Corpus Construction}
\label{app:corpus}

\begin{table*}[t]
  \centering
  \caption{Current Phase-A and external-evaluation rosters. \emph{Cfg} counts materialized sensor
  configurations; \emph{rows} counts stream-window rows before the current quality exclusions and
  subject split, so simultaneous placements are counted separately. $^\dagger$The XRF ear stream was
  acquired at 25\,Hz and stored at 50\,Hz. $^\ddagger$MM-Fit devices have different acquisition
  clocks and are placed on a shared 100\,Hz grid; source rates still bound observability. The seven
  test datasets have not been used for method selection.}
  \label{tab:datasets}
  \scriptsize
  \setlength{\tabcolsep}{3.2pt}
  \begin{tabular}{@{}lrrr@{\hspace{8mm}}lrrr@{}}
    \toprule
    \multicolumn{4}{c}{Phase-A training} & \multicolumn{4}{c}{External evaluation} \\
    \cmidrule(lr){1-4}\cmidrule(lr){5-8}
    Dataset & Hz & Cfg & rows & Dataset (role) & Hz & Cfg & rows \\
    \midrule
    CAPTURE-24~\cite{chan2024capture24} & 100 & 1 & 1,543,573
      & MotionSense~\cite{motionsense} (dev) & 50 & 1 & 4,534 \\
    WISDM~\cite{wisdm} & 20 & 2 & 72,227
      & RealWorld~\cite{realworld} (dev) & 50 & 1 & 11,237 \\
    RealDISP~\cite{realdisp} & 50 & 9 & 63,432
      & Shoaib~\cite{shoaib} (dev) & 50 & 1 & 2,100 \\
    XRF V2~\cite{xrf_v2} & 25/50$^\dagger$ & 6 & 59,442
      & InclusiveHAR~\cite{inclusivehar} (test) & 50 & 1 & 1,260 \\
    DSADS~\cite{dsads} & 25 & 5 & 38,175
      & USC-HAD~\cite{usc_had} (test) & 100 & 1 & 4,350 \\
    NFI-FARED~\cite{nfi_fared} & 100 & 2 & 29,184
      & TNDA-HAR~\cite{tnda_har} (test) & 50 & 1 & 3,353 \\
    PhyTMo~\cite{phytmo} & 100 & 8 & 29,180
      & UT-Complex~\cite{ut_complex} (test) & 50 & 1 & 3,900 \\
    Opportunity~\cite{opportunity_challenge} & 30 & 5 & 28,785
      & MoniPar~\cite{monipar} (test) & 50 & 1 & 12,079 \\
    HARMES~\cite{harmes} & 50 & 1 & 25,583
      & SPAR~\cite{spar} (test) & 50 & 2 & 1,898 \\
    FORTH-TRACE~\cite{forth_trace} & 51.2 & 5 & 12,330
      & Upper Limb Use~\cite{upper_limb_use} (test) & 50 & 4 & 1,124 \\
    UniMiB SHAR~\cite{unimibi} & 50 & 1 & 11,771
      & & & & \\
    HHAR~\cite{hhar} & 50 & 1 & 11,587
      & & & & \\
    KU-HAR~\cite{kuhar} & 100 & 1 & 11,263
      & & & & \\
    UCI-HAR~\cite{ucihar} & 50 & 1 & 10,299
      & & & & \\
    MM-Fit~\cite{mmfit} & 100$^\ddagger$ & 4 & 8,496
      & & & & \\
    SP-SW-HAR~\cite{sp_sw_har} & 100 & 2 & 3,757
      & & & & \\
    PAMAP2~\cite{pamap2} & 100 & 1 & 3,292
      & & & & \\
    MHEALTH~\cite{mhealth} & 50 & 1 & 1,230
      & & & & \\
    \bottomrule
  \end{tabular}
\end{table*}

Phase A uses 18 public datasets and 56 acquisition configurations, whose materialized native-grid
manifests hold 1,963,606 stream-window rows before quality screens and the subject split. The bank
built from the completed encoder contains 250,000 balanced source windows, 546 training subjects,
all 166 canonical training labels, 1,350,834 patch rows, and 2,629,972 sensor rows. Ten datasets are
disjoint from Phase-A training: three form the development roster and seven are sealed for testing.

Table~\ref{tab:datasets} is generated from each stream's \path{grids/native/*/meta.json}. Window
counts are properties of the HALO converters and grid policy, not quantities claimed by the source
dataset papers. The evidence ledger records the aggregation command and distinguishes materialized
rows, quality-screened rows, subject splits, and the balanced 250,000-window bank.

\begin{figure}[ht]
  \centering
  \setlength{\abovecaptionskip}{2pt}
  % =====================================================================
% fig_corpus.tex  --  corpus heterogeneity: placement x sampling rate
%
% DATA PROVENANCE: every point is a real stream in the materialized native
% grids, enumerated 2026-07-26 from training.tokenizer.pretrain_data.CorpusIndex
% (20 training streams over 12 datasets) and data/datasets/*/metadata.json
% (7 held-out datasets).  Bubble area encodes window count; the largest
% (CAPTURE-24 wrist) is 144,142 windows, the smallest (SP-SW-HAR pocket) 1,179.
% =====================================================================
\resizebox{\columnwidth}{!}{%
\begin{tikzpicture}
\begin{axis}[
  width=\columnwidth, height=4.6cm,
  xmin=0.55, xmax=5.55,
  ymin=6, ymax=145,
  ytick={20,50,100},
  yticklabels={20\,Hz,50\,Hz,100\,Hz},
  yticklabel style={font=\scriptsize},
  xtick={1,2,3,4,5},
  xticklabels={waist/hip,pocket,wrist,lower back,head/ear},
  xticklabel style={font=\scriptsize},
  ylabel={sampling rate},
  ylabel style={font=\scriptsize,yshift=-6pt},
  grid=major, grid style={draw=halogray!18},
  axis line style={draw=halogray!60},
  tick style={draw=halogray!60},
  scale only axis,
  clip=false,
]
% ---------- training streams (filled blue discs, diameter ~ #windows^0.4) ----
% diameters below are 2*(1.2 + 4.5*(n/144142)^0.4) pt, n = window count
\tikzset{hdisc/.style={circle,draw=haloblue,fill=haloblue!45,line width=0.4pt,inner sep=0pt}}
\node[hdisc,minimum size=5.44pt] at (axis cs:0.88,50)  {};% hhar        9,587
\node[hdisc,minimum size=5.40pt] at (axis cs:1.00,100) {};% kuhar       9,233
\node[hdisc,minimum size=5.54pt] at (axis cs:1.12,50)  {};% uci_har    10,299
\node[hdisc,minimum size=5.70pt] at (axis cs:1.80,50)  {};% unimib     11,771
\node[hdisc,minimum size=4.88pt] at (axis cs:1.94,50)  {};% xrf l.pkt   5,794
\node[hdisc,minimum size=7.76pt] at (axis cs:2.00,20)  {};% wisdm pkt  39,329
\node[hdisc,minimum size=4.88pt] at (axis cs:2.06,50)  {};% xrf r.pkt   5,794
\node[hdisc,minimum size=3.72pt] at (axis cs:2.00,100) {};% spsw pkt    1,179
\node[hdisc,minimum size=11.4pt] at (axis cs:2.82,100) {};% capture24 144,142
\node[hdisc,minimum size=6.36pt] at (axis cs:2.85,50)  {};% harmes     18,576
\node[hdisc,minimum size=4.36pt] at (axis cs:2.94,100) {};% pamap2      3,186
\node[hdisc,minimum size=3.68pt] at (axis cs:2.95,50)  {};% mhealth     1,104
\node[hdisc,minimum size=7.26pt] at (axis cs:3.00,20)  {};% wisdm wr   30,876
\node[hdisc,minimum size=4.88pt] at (axis cs:3.05,50)  {};% xrf l.wr    5,794
\node[hdisc,minimum size=5.86pt] at (axis cs:3.06,100) {};% nfi wrist  13,260
\node[hdisc,minimum size=4.88pt] at (axis cs:3.15,50)  {};% xrf r.wr    5,794
\node[hdisc,minimum size=4.20pt] at (axis cs:3.18,100) {};% spsw wrist  2,578
\node[hdisc,minimum size=5.86pt] at (axis cs:4.00,100) {};% nfi back   13,260
\node[hdisc,minimum size=4.88pt] at (axis cs:4.92,50)  {};% xrf glasses 5,794
\node[hdisc,minimum size=4.88pt] at (axis cs:5.08,50)  {};% xrf earbuds 5,794
% ---------- held-out datasets (open orange triangles) ----------
\addplot[only marks,mark=triangle*,mark size=2.6pt,color=haloorange,
         fill=halolight2,line width=0.7pt]
  coordinates {
    (1.26,50) (1.38,50) (1.26,100)
    (2.20,50) (2.32,50)
    (3.30,50) (3.42,50)
  };
% ---------- annotations ----------
\node[font=\scriptsize,color=haloblue!75!black,anchor=west] at (axis cs:2.95,128)
  {\textbf{CAPTURE-24} 144k windows};
\draw[haloblue!70,thin] (axis cs:2.86,106) -- (axis cs:2.93,124);
\node[font=\scriptsize,color=halogray,anchor=west,align=left] at (axis cs:4.05,24)
  {smart glasses\\ \& earbuds};
\draw[halogray,thin] (axis cs:5.00,44) -- (axis cs:4.90,31);
\node[font=\scriptsize,anchor=west,align=left] at (axis cs:0.62,13)
  {\textcolor{haloblue!75!black}{$\bullet$ 20 training streams}\;\;
   \textcolor{haloorange!85!black}{$\blacktriangle$ 7 held-out}};
\end{axis}
\end{tikzpicture}}
\caption{\textbf{The corpus does not cover a grid; it covers a scatter.} Each blue disc is
one of the 20 sensor streams HALO trains on, positioned by body placement and native
sampling rate, with area proportional to the number of $6$\,s windows it contributes; orange
triangles are the 7 held-out evaluation datasets. Rates span a $5\times$ range and placements
include two---smart glasses and wireless earbuds---that no established HAR benchmark
contains. Held-out cells are not simply new subjects: they combine placements and rates in
proportions the training corpus never sees. Absorbing this by resampling to a common rate
would discard the content above $10$\,Hz that only the $50$/$100$\,Hz streams can
resolve.}
\label{fig:corpus}

\end{figure}

The expanded training roster deliberately includes body-worn stress configurations from DSADS and
Opportunity; an earlier draft incorrectly listed both as excluded. Six newly added training
sources contribute most of the placement coverage: DSADS, FORTH-TRACE, Opportunity, RealDISP,
MM-Fit, and PhyTMo. MoniPar, SPAR, and Upper Limb Use are held out.

\textbf{Sources not used in the named recipe.} HAPT~\cite{hapt} reuses the UCI-HAR participants and
recordings, so it is excluded to avoid near-duplicate subjects. RecGym~\cite{recgym} distributes
per-axis min--max-normalized data that no longer preserve the physical scale used by the tokenizer.
The MobiAct~\cite{mobiact} download did not produce a materialized grid and is not counted.
HARTH~\cite{harth} is converted but has no materialized grid. ExtraSensory, NHANES, HMOG, and
KneE-PAD are buildable optional scale and stress sources that must be requested explicitly, and are
not members of the reported 18-dataset run.

\section{Baseline Comparison}
\label{app:baselines}

\pendingsection{No baseline has been run against the completed encoder. Models receive different
input information, so any future comparison must disclose both weight source and input contract:
LiMU-BERT and CrossHAR require corpus-matched retraining, while released harnet, UniMTS and NormWear
weights were trained on different corpora. ConSE~\cite{conse} can translate a fixed training
classifier into candidate-text scores, but it does not make the underlying encoder open-label or
supply an enrollment mechanism. Results will be reported only when every row shares the support and
query cohorts used for HALO.}

\section{Tokenizer Validation}
\label{app:tokenizer}

\pendingsection{The implementation carries unit guards for DFT overflow, acquisition-rate
observability, sensor folding, and axis validity, and those tests pass. They establish that the code
computes what it claims, not that the resulting features are empirically rate-equivalent or that the
signed low-frequency feature is useful. Both are open measurements. Their generating script must
record input windows, anti-alias filtering, source and stored rates, checkpoint hash, random seed,
and raw feature arrays, and the signed low-frequency ablation must use the same evaluation cohort as
the main result and report uncertainty.}

\section{Phase-A Specification}
\label{app:phasea}

The reported run uses a fixed filterbank, sensor-granularity tokens, width 256, six layers, eight
heads, batch size 256, 30,000 steps, learning rate $3\times10^{-4}$, 1,000 warmup steps, weight decay
$0.05$, gradient clipping at $1.0$, and seed 20260718. The two top-level loss coefficients are fixed
after calibration at step 2,000: 7.053501 for the JEPA family and 0.653232 for VICReg. The teacher
decay is $0.996$. Descriptor retrieval carries coefficient zero and is not trained.

Patches are one second on a single resolution. A dataset with $n$ windows is sampled proportional to
$n^{0.25}$ with a maximum share of $25\%$; subjects are sampled proportional to the square root of
their window count. Activity labels do not enter the training sampler or loss. The selected
checkpoint is step 27,000, which maximizes the predefined subject-disjoint label--stream-macro
validation score over the completed run.

\section{Evidence Artifact and Controls}
\label{app:phaseb}

The active bank has 250,000 source windows, 56 configurations, and all 166 canonical labels. It
records 1,350,834 valid patch rows and 2,629,972 per-sensor rows, and it is bound to the completed
encoder checkpoint and to the sentence-embedding path used to build it. A gate artifact must
additionally record the resolvability table and bank it was fitted against. Evaluation refuses
legacy banks or foreign artifacts unless the mismatch is explicitly acknowledged.

The prediction budget is 64 evidence rows per query patch and candidate, sensor-bias similarity
carries weight zero, and for semantically unfamiliar sensor or concept text the gate blends toward a
neutral value.

\pendingsection{No gate has been fitted against this bank, so none of the telemetry below has been
produced. When a gate is reported, four items are mandatory rather than optional: gate mean, spread,
extrema and vetoed-row fraction; same-dataset retrieval with the acquisition-bias blend both off and
on; held-out concept, dataset and body-region errors relative to constant and per-concept baselines;
and within-concept correlation between the gate latents and the train-only resolvability table.
Every enrollment result must share one deterministic query cohort and nested support prefixes across
the gate, all-one admissibility, support removal, support-label shuffling, prototype, and ridge
controls.}

\section{Evidence Ledger}
\label{app:ledger}

The paper source includes \texttt{EVIDENCE\_LEDGER.md}. Each quantitative manuscript statement is
classified as verified, development-only, sealed, or pending. A verified entry names the producing
command or log, source-data manifest, seed, checkpoint, and artifact fingerprint. Values without
that chain are not printed as results. This policy is also why the earlier geometry, gravity,
rate-equivalence, DFT-cosine, deployment, and phone-latency numbers are absent from this revision.
